\begin{russian}
\thispagestyle{empty}
\begin{center}
{\Large \textbf{АННОТАЦИЯ}\par}
\end{center}
\vspace{0.5cm}

Современные детекторы аномалий для АСУ ТП показывают высокие значения F1-меры на наборах данных, на которых они обучались, однако их рейтинги существенно меняются при применении временно-чувствительных метрик оценки, и они теряют обобщающую способность при переносе между испытательными стендами --- оба эффекта скрыты за заголовочными числами, доминирующими в литературе.  Настоящая дипломная работа сравнивает шесть детекторов (Isolation Forest, One-Class SVM, плотный автокодировщик, LSTM-автокодировщик, USAD, TranAD) на наборах данных HAI 21.03 и Morris (магистральный газопровод) с использованием трёх временно-чувствительных метрик, проводит исследование межстендового переноса, разделяющее перенос представлений и перенос порога принятия решений, а также оценивает посенсорную атрибуцию на уровне технологических процессов.

Три ключевых вывода работы: (1)~межстендовый перенос с калибровкой порога на исходном домене вырождает каждый детектор в тривиальный классификатор «все --- атаки», соответствующий априорной вероятности класса в целевом домене, тогда как калибровка на целевом домене вскрывает осмысленное разделение классических и современных моделей, невидимое в стандартном протоколе; (2)~эксперимент «исключение одного процесса» внутри HAI разделяет сдвиг распределения признаков и сдвиг априорной вероятности класса, демонстрируя адресную «слепоту» обнаружения при удалении специфических для процесса сенсоров; (3)~плотный автокодировщик --- слабейший детектор по F1 --- оказывается единственной моделью, чья попризнаковая атрибуция стабильно превосходит случайный базовый уровень для всех трёх атакуемых процессов HAI, что устанавливает компромисс между обнаружением и локализацией.

Вкладом работы являются: воспроизводимый протокол сравнительного тестирования, открытая реализация моделей в рамках унифицированного интерфейса детектора, а также методологическая рекомендация сообществу обнаружения аномалий в АСУ ТП: представлять полную таблицу метрик, режимов калибровки и результатов атрибуции вместо одного заголовочного значения F1.  Каждое число в документе воспроизводимо из зафиксированных в репозитории YAML-конфигураций.

\vspace{0.5cm}
\noindent\textbf{Ключевые слова:} промышленные системы управления; обнаружение аномалий; кибер-физическая безопасность; межстендовый перенос; временно-чувствительные метрики; атрибуция; HAI; Morris.
\end{russian}
\newpage
